\chapter*{Abstract}

Speech Enhancement (SE) has undergone significant advancements in recent years, moving from traditional algorithmic methods to more advanced deep learning-based models. These deep learning models have greatly improved speech quality and intelligibility.  A major development in SE is personalized speech enhancement (PSE),  where models are trained to enhance the speech of specific speakers. This development was evident in the DNS Challenge 2023, where the top-performing models were personalized \cite{dubey2023}. Our goal is to further enhance the challenge by introducing real-time constraints, which will impose additional limitations on the model's architecture and size.

In this work, we explore various models and architectures in SE. We propose two models: the Personalized Fast FullSubNet (PFFSN) and the Personalized Denoiser (PDenoiser). Both models build on established SE architectures. We extend these models by adding personalization support. PFFSN, an adaptation of FFSN \cite{hao2023fast}, uses a subband approach to classify and amplify the subbands where the targeted speaker's voice resides. PDenoiser, an adaptation of Denoiser \cite{defossez2020}, suppresses both stationary and non-stationary noises as well as non-primary speakers, all while operating in the time domain. Our experiments demonstrate that PFFSN can compete with state-of-the-art models in terms of speech quality and intelligibility, as evidenced by its performance in the DNS Challenge.

% PFFSN is a lightweight model that uses a Fast Fourier Transform (FFT) to extract features from the input speech signal. The extracted features are then passed to a small neural network that enhances the speech signal. PDenoiser is a personalized speech enhancement model that consists of two major modules: a speaker embedding module and a denoising module. The speaker embedding module extracts speaker-specific features from the input speech signal, while the denoising module enhances the speech signal based on the extracted speaker features.

We also introduce a novel Arabic dataset specifically designed for personalized speech enhancement, which was meticulously collected independently. This dataset comprises 30 speakers, with a demographic distribution of 73\% male and 27\% female participants. The training-validation set includes 30,000 audio clips amounting to a total of 83 hours of recorded speech, while the test set encompasses 1,200 audio clips totaling 3.3 hours.

Finally, we are pleased to introduce a real-time desktop application that seamlessly integrates trained models. This application features a user-friendly graphical interface (GUI) for effortless operation and supports on-device real-time audio processing. It incorporates optimizations for streaming procedures, including adjustable frame size and multi-threading. Impressively, it achieves a low real-time factor (RTF) of 0.21.